\documentclass[12pt]{extarticle}
\def\activityid{F03-FAC-v3}\usepackage[letterpaper,margin=.65in,top=.60in,bottom=.65in]{geometry}
\usepackage{fix-cm}
\usepackage[T1]{fontenc}\usepackage[default]{sourcesanspro}
\usepackage{microtype,tikz,array,tabularx,fancyhdr,amsmath,amssymb}
\usepackage[hidelinks]{hyperref}\usetikzlibrary{calc,arrows.meta}
\setlength{\parindent}{0pt}\setlength{\parskip}{7pt}\setlength{\headheight}{0pt}\setlength{\footskip}{26pt}
\setlength{\emergencystretch}{2em}\pagestyle{fancy}\fancyhf{}\renewcommand{\headrulewidth}{0pt}\renewcommand{\footrulewidth}{.4pt}
\fancyfoot[L]{\scriptsize Bellingham Math Circle / Week 3 / \activityid}\fancyfoot[R]{\scriptsize \thepage}
\definecolor{ink}{gray}{.12}\definecolor{soft}{gray}{.95}\color{ink}
\newcommand{\PageTitle}[2]{{\fontsize{10}{12}\selectfont\bfseries WEEK 3\quad CODE MACHINES\hfill #1\par}\vspace{3pt}{\fontsize{26}{29}\selectfont\bfseries #2\par}\vspace{2pt}}
\newcommand{\NameLine}{{\small Name: \rule{2.65in}{.4pt}\hfill Date: \rule{1.35in}{.4pt}\par}\vspace{4pt}}
\newcommand{\Task}[2]{\par\vspace{3pt}{\large\bfseries #1}\par #2\par}
\newcommand{\Subhead}[1]{\par\vspace{3pt}{\large\bfseries #1}\par}
\newcommand{\RuleBox}[1]{\par\begingroup\setlength{\fboxsep}{8pt}\colorbox{soft}{\parbox{\dimexpr\linewidth-16pt\relax}{#1}}\endgroup\par}
\newcommand{\WriteBox}[2]{\par\textbf{#1}\par\vspace{2pt}\begin{tikzpicture}\draw[black!40,line width=.5pt] (0,0) rectangle (\dimexpr\linewidth-.5pt\relax,#2);\end{tikzpicture}\par}
\newcommand{\StudentFooter}{\vfill{\small Try it with objects. Explain by pointing, talking, or drawing.}\par}
\newcommand{\LampRing}[3]{\begin{tikzpicture}[x=1in,y=1in]
\foreach \i in {1,...,#1}{\coordinate (v\i) at ({90-360*(\i-1)/#1}:#2);}
\foreach \i in {1,...,#1}{\pgfmathtruncatemacro{\j}{mod(\i,#1)+1}\draw[thick] (v\i)--node[midway,fill=white,inner sep=2pt,font=\small]{\i\j}(v\j);}
\foreach \i in {1,...,#1}{\node[circle,draw,fill=white,minimum size=.52in,inner sep=0pt,font=\bfseries] at(v\i){\i};}
\foreach \i in {#3}{\node[circle,draw,fill=ink,text=white,minimum size=.52in,inner sep=0pt,font=\bfseries] at(v\i){\i};}
\end{tikzpicture}}
\newcommand{\Lights}[1]{\begin{tikzpicture}[x=.55in,y=.55in,baseline=-.10in]\foreach \v [count=\i] in {#1}{\ifnum\v=1\fill (\i,0) circle(.16in);\else\draw (\i,0) circle(.16in);\fi}\end{tikzpicture}}
\newcommand{\ClockRing}[2]{\begin{tikzpicture}[x=1in,y=1in]\draw[black!30] (0,0) circle(#2);\pgfmathtruncatemacro{\n}{#1-1}\foreach\i in {0,...,\n}{\node[circle,draw,fill=white,minimum size=.35in,inner sep=0pt] at({90-360*\i/#1}:#2){\i};}\draw[-{Stealth},thick] (-.35,.15) arc(155:25:.38);\node[font=\small] at(0,-.18){clockwise};\end{tikzpicture}}
\newcommand{\Key}[2]{\begin{center}\renewcommand{\arraystretch}{1.55}\begin{tabular}{c|*{#1}{c}}in&#2\end{tabular}\end{center}}


% v3 student pages use one running header and numbered tasks only.
\newcommand{\StudentHeader}[1]{{\fontsize{10}{12}\selectfont\bfseries Week 3 / Code machines / #1\par}\vspace{10pt}}
\newcommand{\Problem}[2]{\par\vspace{4pt}\textbf{Problem #1:} #2\par}
\newcommand{\Space}[1]{\begin{tikzpicture}\draw[black!35] (0,0) rectangle (\dimexpr\linewidth-.5pt\relax,#1);\end{tikzpicture}\par}
\newcommand{\Record}[2]{#1\par\Space{#2}}

\begin{document}

\PageTitle{FACILITATOR / 1}{Prepare the code machines}
\RuleBox{v3 revision, September 27, 2026: unpiloted. Week 1's reported need for more concrete work and clearer sequencing is the evidence for this redesign; it is not evidence that these code tasks have been tested.}
\Subhead{Materials and prerequisites}
Give K children two sets of circle/triangle/square slips, middle children A--D cards, and upper children A--E (F--I in reserve). Bring 13 counters, paper and pencils. Print first pages for three K, three middle, and one upper child (seven starting sheets). Keep continuation masters ready; copy later pages when chosen. The extra can remain digital. Read aloud and scribe freely.\par
\textbf{K:} recognize shapes and count to three. \textbf{Middle:} recognize letters as labels, count to six, and follow a fixed mapping. \textbf{Upper:} track five positions and coordinate return times; multiples can be generated by moving counters. \textbf{Extra:} common multiples to 20 and organizing splits by largest part. Factorials and formal group notation are not prerequisites.
\Subhead{Handle, then launch: 0--15 minutes}
Let children arrange cards into messages during 0--10 minutes. At 10--15 minutes, gather around the circle-to-triangle-to-square-to-circle key. Point to one input, follow its arrow, and place its output below it. Encode a three-symbol message one position at a time, keeping the original row until the new row is complete. Have a child read the output and check it against the arrows. Explain that the next turn uses the new row as its input. Then distribute starting tasks. Do not introduce loop classification yet.
\Subhead{A flexible hour}
\textbf{15--35:} K sends and reverses messages; middle experiments with messages then P/Q; upper compares five-loop and 3+2 keys. Adult A stays with K; organizer checks complete-row substitution. Compare a first return when useful. \textbf{35--40:} stand, stretch, reset. \textbf{40--55:} repeat, draw loops, or choose an extension when ready. \textbf{55--60:} share a key and its return, then tidy. Give one page at a time.
\Subhead{Stops and representation gates}
\textbf{K:} sending and decoding several messages is a complete stopping point. After repeated single-symbol turns, introduce the arrow loop as a record of those moves. The swap/fixed key comes later; the broken key is optional.\par
\textbf{Middle:} stop after P/Q and a four-turn construction. Demonstrate how one row of the key becomes an arrow before drawing all loops. Let children build each size arrangement before asking whether the cases are complete.\par
\textbf{Upper:} stop after the 5 versus 3+2 comparison. Move to simultaneous-return markings, then classification (6--7). Composition (8--11) is a separate direction. \textbf{Extra:} stop after constructing the 15-turn key; classification is the route to its upper bound, nine letters a later reserve.
\Subhead{Proof conversations}
Ask ``Is every letter home?'' ``Can two inputs point to one output?'' ``What can be left after the largest loop?'' Keep the loop theorem, exhaustive bounds, and commuting explanations in conversation; the full arguments follow. Do not read a supplied result as if the child proved it.
\newpage
\PageTitle{FACILITATOR / 2}{Reversibility and loops}
\Subhead{K--1 checked results}
The three-shape key sends circle--triangle--circle to triangle--square--triangle. The input giving square is triangle. Every symbol returns after three turns; the order of names or the starting point does not affect that count. Under the changed key, circle and triangle swap while square stays put. All are home after two turns; square is also home after one. A key sending both circle and triangle to square loses information: the receiver cannot determine which was sent.
\Subhead{Middle Problems 1--5}
P sends ABBA to \textbf{BCCB}; decoding CAAC gives \textbf{BCCB}. A one-to-one key preserves equality and inequality between positions, so ABBA cannot become ABCD. Its first and last outputs must match, as must the middle pair.
\begin{center}\renewcommand{\arraystretch}{1.4}\begin{tabular}{c|cc}Turn&P&Q\\\hline 0&ABCD&ABCD\\1&BCAD&BADC\\2&CABD&ABCD\\3&ABCD&BADC\\4&BCAD&ABCD\end{tabular}\end{center}
P has a three-cycle A--B--C--A and fixed D, so order 3. Q has two swaps A--B and C--D, so order 2. A single fixed letter does not certify that the whole key has returned.
\Subhead{Middle Problems 6--8: complete classification}
A four-cycle A--B--C--D--A has order 4. The possible cycle-size partitions and their orders are:
\begin{center}\begin{tabular}{c|ccccc}Loop sizes&4&3+1&2+2&2+1+1&1+1+1+1\\\hline Order&4&3&2&2&1\end{tabular}\end{center}
These exhaust four letters by sorting on the largest loop. Order 6 is impossible with four letters; separate loops of size 3 and 2 require five. Counting all possible keys gives \(4\cdot3\cdot2\cdot1=24\), but the classification uses five cases. Do not replace the investigation with a factorial exercise.
\Subhead{Why every reversible key splits into loops}
Follow arrows from any symbol. Finiteness forces a repeated symbol. The first repeat must be the starting symbol: if a later point were first repeated, it would have two distinct predecessors, violating the one-to-one rule. Remove that loop and repeat with an unused symbol. Loops are disjoint and exhaust all symbols. On a loop of length a, the first return is a and all returns occur at multiples of a. All loops return together at their least common multiple. This is a child-accessible proof, not just a pattern from examples.
\newpage
\PageTitle{FACILITATOR / 3}{A maximum and noncommuting keys}
\Subhead{Upper: the five-letter example}
Successive messages are ABCDE, BCAED, CABDE, ABCED, BCADE, CABED, ABCDE. The 3-cycle returns on turns 3,6,9,\ldots; the 2-cycle on 2,4,6,\ldots. The first common return is \textbf{6}. A five-cycle takes only 5, so the longest individual loop need not be optimal.
\begin{center}\renewcommand{\arraystretch}{1.3}\begin{tabular}{c|rrrrrrr}Sizes&5&4+1&3+2&3+1+1&2+2+1&2+1+1+1&1+1+1+1+1\\\hline Order&5&4&6&3&2&2&1\end{tabular}\end{center}
Let children organize the blank table before showing the complete list below. Largest loop 5 or 4 leaves just one split; largest 3 leaves 2 or 1+1; largest 2 leaves 2+1 or 1+1+1; largest 1 forces five singletons. Ask why these exhaust the leftovers. This complete list plus the loop theorem proves no five-letter key exceeds 6; the displayed key attains it. The mathematical work is construction \emph{and} upper bound. If reading the full table is a barrier, arrange the five letter cards into loops and enumerate the same seven cases physically.
\Subhead{Upper: P then Q is not Q then P}
P swaps A/B; Q swaps B/C; both fix D/E. P then Q takes ABCDE to BACDE to \textbf{CABDE}. Q then P takes ABCDE to ACBDE to \textbf{BCADE}. At letter A alone: A goes to B then C in the first order, but A then B in the second. Thus the maps differ.
The composite P then Q is the loop A--C--B--A with D/E fixed: its order is \textbf{3}, even though each factor has order 2. To undo it, apply Q then P. Disjoint swaps, such as A/B and D/E, commute: neither changes any input moved by the other. Their common composite has two 2-cycles and order 2.
\Subhead{Extra: the eight-letter optimum is 15}
Use a 5-cycle on A--E and a 3-cycle on F--H. Their first common return is 15. Exhaustive upper bound by largest loop: 8 gives 8; 7 leaves one and gives 7; 6 leaves at most a 2-cycle and gives 6; 5 leaves three letters, giving at most \(\mathrm{lcm}(5,3)=15\); 4 leaves four, with largest possible LCM 12 from 4+3+1; if all loops have size at most 3 their LCM divides 6. Thus nothing exceeds 15.
For nine letters, 5+4 gives \textbf{20}. For completeness: largest 9,8,7,6 give maxima 9,8,14,6; largest 5 leaves at most four and gives at most 20; largest 4 gives at most 12; all at most 3 gives at most 6. Adding a letter never decreases the maximum (fix the new letter) but need not strictly increase it: sizes 5 and 6 both have maximum 6. On six letters, partitions with largest 6,5,4,3,2,1 have maximal orders 6,5,4,6,2,1.
\newpage
\PageTitle{FACILITATOR / 4}{The exact research question}
\Subhead{Undergraduate mathematics}
A reversible substitution on n symbols is a permutation, an element of the symmetric group. A repeated key is a power of this permutation; its return time is its order. Disjoint cycles, least common multiples, inverses, and noncommuting products are standard abstract algebra. No formal group axioms are needed in the child-facing activity. The words ``key'' and ``code'' describe the entry point; these small deterministic substitutions are not modern secure encryption.
\textbf{Thomas Judson}, \emph{Abstract Algebra: Theory and Applications}, Chapter 5, \S5.1 ``Definitions and Notation,'' especially disjoint-cycle examples and Theorem 5.9 on disjoint-cycle decomposition. See \href{https://people.hsc.edu/faculty-staff/blins/books/JudsonAbstractAlgebra.pdf}{university-hosted textbook}, printed pp. 59--63. Our explicit loop proof supports the activity independently of theorem numbering.
\Subhead{A direct research lineage}
\textbf{Marc Del\'eglise, Jean-Louis Nicolas, and Paul Zimmermann}, \emph{Landau's function for one million billions} (2008 arXiv manuscript), \S1, defines g(n), the maximum order of a permutation of n symbols, and relates it to least common multiples of partitions. The paper develops algorithms for very large n: \href{https://arxiv.org/abs/0803.2160}{arXiv:0803.2160}; \href{https://arxiv.org/pdf/0803.2160}{paper, introduction pp. 1--3}.
The classroom's g(5)=6 and g(8)=15 are \emph{exact small instances of the same optimization problem}. We do not claim that finding these small values is new research. Efficient computation and large-n growth are the research directions. The extra packet asks for a proved maximum, not a decorated encoding exercise.
\Subhead{Downloaded sources and pedagogical choices}
Rozhkovskaya, \emph{Math Circles for Elementary School Students}, Lesson 2, problems 2.2--2.8 (EPUB section 12), uses secret codes and shifts. We retain the physical key and short-message entry while replacing the core with arbitrary reversible substitutions and loop structure. This week's keys need not be shifts; week 4 imposes a fixed circular order.
\emph{Math Circle by the Bay}, preface pp. viii--x, supports a deep theme accessible across levels, manipulatives, explanation, and spare challenges. Rozhkovskaya Lesson 3, ``At the lesson'' item 1, supports attempted solutions and explanations before tabular organization. Our exact keys, tables, staffing, and timing are new adaptations.
\Subhead{Prior use and verification}
Lowell Handout 8, problem 8.1, already used divisibility for distributing apples among 3,4,5,6 guests. The LCM idea here is revisited through synchronized permutation loops, with new questions about maximal orders. No saved use log proves these particular machines were taught. Record actual instances as F03-K/M/U/X-v3 and note which keys or partitions children saw.
\texttt{verify.py} checks every permutation for n=4,5,6,8,9 against the claimed maxima and verifies all message traces and compositions. It is a separate check of the finite results, not a substitute for the loop and partition proofs.

\newpage
\PageTitle{FACILITATOR / 5}{v3 checks and observations}
\Subhead{Added concrete instances}
\textbf{K 2:} circle--triangle--circle becomes triangle--square--triangle; square--square--triangle becomes circle--circle--square; triangle--circle--square becomes square--triangle--circle. \textbf{3:} received triangle--circle--triangle decodes circle--square--circle. \textbf{4--6:} each start returns at turn 3; the whole row just before return is square--circle--triangle. \textbf{7--8:} swapped shapes return in 2, the fixed square at every turn; whole return is 2 versus 3 for the original key. \textbf{9:} either circle or triangle could have produced square.\par
\textbf{Middle 1:} P gives BCCB, BDBD, BCAD. \textbf{2:} CAAC decodes BCCB; DBDB decodes DADA. \textbf{3:} every key preserves the equal-position pattern of ABBA. A child may test many keys; the proof is that identical inputs share an output and distinct inputs have distinct outputs. \textbf{4--5:} P/Q rows and loops are on page 2. \textbf{6--8:} four-letter loop returns at 4; 3+1 at 3; 2+2 and 2+1+1 at 2; four fixed letters at 1. These five partitions rule out 6; five letters with loops 3+2 achieve 6.
\Subhead{Upper: experiments before classification}
\textbf{1:} R rows are ABCDE, BCDEA, CDEAB, DEABC, EABCD, ABCDE; S rows are on page 3. \textbf{3:} 4+1 returns at 4; 2+2+1 at 2; 3+1+1 at 3. \textbf{4:} size 3 marks 3,6,9,12; size 2 marks 2,4,6,8,10,12. First common return is 6. Introduce ``least common multiple'' as a name for this record, not an entry requirement. \textbf{6--7:} the seven partitions and completeness proof are on page 3. Do not disclose the count before children organize cases. \textbf{8--11:} page 3 gives both compositions and inverse; P/T both orders give BACED.
\Subhead{Extra: bounds after constructions}
\textbf{1--2:} returns for 8,4+4,6+2,5+3 are 8,4,6,15. Two loops of equal length return together at that length, not its square. \textbf{3--4:} for largest loops 8,7,6,5,4,at most 3, row bounds are 8,7,6,15,12,6. In the last row, every loop has size 1,2,3, so turn 6 returns all loops; no list of ten partitions is necessary. \textbf{5:} nine-letter examples 5+4,7+2,4+3+2 return at 20,14,12. \textbf{6:} the nine-letter bounds are on page 3. \textbf{7:} six-letter trials return 6,6,5. The six-letter partition bound on page 3 proves the plateau at g(5)=g(6)=6; examples alone do not prove a maximum.
\Subhead{Observe, then adapt}
Record actual keys/problems used, whether every letter was changed from the old row, a child's return prediction, success reading/drawing loops, copying that displaced exploration, and requests to continue. Keep observations separate from explanations of why they happened and next-time proposals. Mark proof status explicitly. All v3 expansions are unpiloted design proposals.

\end{document}
